\section{Methods}
\label{sec:methods}

\subsection[Datasets and signal processing]{\chg{Datasets and signal processing}}
\label{sec:datasets}

\chg{The primary analysis used GRABMyo v1.1.0 \citep{pradhan2022multiday,jiang2024grabmyo}, distributed through PhysioNet \citep{goldberger2000physionet}, comprising recordings from \GRABMyoParticipants\ able-bodied participants collected over \GRABMyoSessions\ sessions on different days.} Participants were \GRABMyoMale\ male and \GRABMyoFemale\ female (age \GRABMyoAgeMean\ $\pm$ \GRABMyoAgeSD\ years). Signals were recorded with an EMGUSB2+ amplifier (OT Bioelettronica) at 2048~Hz; in each session, participants performed each gesture in 7 trials of 5~s. \chg{Sixteen monopolar electrodes were positioned in two circumferential rings of eight electrodes around the forearm (F1--F8 and F9--F16), separated longitudinally by 2~cm,} with electrode $j$ of one ring aligned with electrode $j+8$ of the other. Mean forearm circumference at the ring position was \GRABMyoCircMean\ $\pm$ \GRABMyoCircSD~cm, giving about \GRABMyoArcSpacing~cm between neighbouring electrodes in a ring. \chg{Participants performed 16 active hand and wrist gestures together with rest; the primary recognition task considered the 16 active gestures only,} for which chance Macro-F1 is about 0.063. All \GRABMyoFilesVerified\ distributed files matched the dataset's published SHA-256 checksums, and all \GRABMyoRecords\ records were analysed.

Terminology is used as follows. An \emph{electrode} is a physical recording site and a \emph{channel} is the signal it provides to the classifier; in the monopolar analyses each electrode yields one channel, so selecting electrodes and selecting channels are the same operation. The \emph{montage} is the spatial arrangement of the electrodes and the \emph{array} is the complete set. A \emph{bipolar pair} is the difference between two electrodes, derived from the monopolar recordings in one exploratory analysis. The subset size $k$ always counts the selection unit of the analysis in question.

\chg{NinaPro DB5 \citep{pizzolato2017comparison} was analysed separately as an exploratory replication dataset and was never pooled with GRABMyo.} It provides \DBFiveSubjects\ able-bodied subjects (\DBFiveMale\ male, \DBFiveFemale\ female; age \DBFiveAgeMean\ $\pm$ \DBFiveAgeSD\ years) recorded with two Myo armbands (16 channels, 200~Hz). We used exercise~1 (12 finger movements, 6 repetitions each), segmented with the refined movement and repetition labels supplied with the dataset. The DB5 values are 8-bit integer codes, not physical units, and its channel numbers have no anatomical correspondence with those of GRABMyo, so only the behaviour of the procedure is compared across datasets.

\paragraph{Ethics statement} This study is a secondary analysis of previously published, de-identified datasets and involved no new participant recruitment or data collection. The original data collections were conducted in accordance with the Declaration of Helsinki and approved by the Office of Research Ethics of the University of Waterloo (ORE\#~31346; GRABMyo) and the Ethics Commission of the Canton of Valais, Switzerland (NinaPro DB5); informed-consent procedures are described in the dataset publications \citep{pradhan2022multiday,pizzolato2017comparison}. [AUTHOR INPUT REQUIRED: state whether an institutional review board determination or exemption was obtained for the present secondary analysis.]

\paragraph{Sex as a reported variable} Sex was taken as reported by the source datasets. Sex-stratified analyses were not prespecified, and the study was not designed to estimate sex differences in recognition performance or electrode ranking; no sex-specific inference is made.

\chg{GRABMyo signals were band-pass filtered between 20 and 450~Hz using a fourth-order zero-phase Butterworth filter and notch filtered at 50~Hz} ($Q = 30$), the mains frequency identified from the spectra (mean band-power ratio 10.35 at 50~Hz and 0.62 at 60~Hz). Because the filters are zero-phase, the processing is non-causal and suitable only for offline evaluation. \chg{Signals were segmented into 250-ms windows with 50\% overlap, without allowing windows to cross trial boundaries. Ten conventional time- and frequency-domain features were extracted from each electrode: mean absolute value, root mean square, variance, waveform length, zero crossings (ZC), slope sign changes (SSC), integrated EMG, Willison amplitude (WAMP), mean frequency and median frequency. The resulting representation contained 160 features per window} (\GRABMyoWindows\ windows in total, \GRABMyoActiveWindows\ from active gestures). \chg{Signal-processing parameters were fixed consistently across the analysis.} The amplitude threshold of the ZC, SSC and WAMP features was set by one rule for every representation: 1.15 times the median absolute first difference of the signal, rounded to one significant figure. For GRABMyo it was computed once, from 54 unfiltered recordings of 6 participants, giving $10^{-2}$~mV; it was therefore not estimated within each outer fold. \chg{A leave-one-out check} gave the same value in every fold (supplementary methods~S2), so fold-wise estimation would have produced identical features. DB5 was band-pass filtered at 10--95~Hz without a notch filter, with 250~ms windows (50 samples) and a threshold of 7 quantisation steps from the same rule; the leave-one-out recomputation gave the same threshold in every DB5 fold.

\subsection[Subject-independent electrode ranking and model evaluation]{\chg{Subject-independent electrode ranking and model evaluation}}
\label{sec:ranking}

\chg{All primary evaluations followed an outer leave-one-subject-out (LOSO) design (figure~\ref{fig:workflow}). For each held-out participant $p$, all preprocessing statistics, model fitting and electrode ranking were derived exclusively from the remaining participants $R_{-p}$. The held-out participant contributed only to the final performance estimate.}

\begin{figure}[t]
\centering
\includegraphics[width=\textwidth]{figures/fig2_workflow.pdf}
\caption{Study workflow. For each held-out participant $p$, feature standardisation, model fitting and electrode ranking use only the remaining participants $R_{-p}$. Attributions are computed on participant-disjoint inner folds of $R_{-p}$, and models restricted to the top-ranked $k$ electrodes are retrained on $R_{-p}$ and tested on $p$. The same workflow was applied separately to each dataset.}
\label{fig:workflow}
\end{figure}

\chg{XGBoost was used as the primary classifier for the electrode-selection experiments because the analysis required repeated model retraining across participants, electrode counts and selection strategies.} It was designated in the analysis plan on these computational grounds; the plan was written before the full-cohort classifier comparison had finished, although results from an earlier, superseded run on 15 participants were available. XGBoost used 300 trees, depth 6, learning rate 0.1, row and column subsampling of 0.8, histogram-based tree construction and a multiclass softmax objective, and was trained on all \FoldTrainingWindows\ training windows of each fold. \chg{Support vector machine and random forest models were additionally evaluated using the complete array to provide classifier-level context:} a support vector machine with a radial basis function kernel ($C = 1$, kernel width set from the feature variance), trained on a stratified sample of 40,000 windows, and a random forest (200 trees, minimum leaf size 2) trained on 200,000 windows. These sample sizes were set by computational cost and chosen on two folds (supplementary methods~S1). Halving the support vector machine's training sample on one fold lowered Macro-F1 from \SVMCalibFortyK\ to \SVMCalibTwentyK, so its result is likely a lower bound. Classifier hyperparameters were fixed for the reported runs. Because this earlier budget calibration used test-fold outcomes, the development process was not fully isolated from the evaluation cohort.

\chg{Electrode importance was estimated within each outer training set using $K = 5$ participant-disjoint inner folds. For each inner fold $j$, an XGBoost model was trained on the remaining inner participants and TreeSHAP values were evaluated on the participants excluded from that model.} Each inner model used a random sample of at most 200,000 windows standardised with inner-training statistics only, and exact tree-path-dependent TreeSHAP attributions on raw model margins \citep{lundberg2020local} were computed for up to 20,000 windows sampled at random from the inner-held-out participants. Because participants contribute similar numbers of windows (median \TestWindowsPerParticipant), each participant carries roughly equal weight. In every fold we verified that no inner-held-out participant contributed to the corresponding inner training set. \chg{Absolute attributions were averaged over the $G = 16$ classes and the $F = 10$ features belonging to each electrode,}
\begin{equation}
\Phi_{i,c} = \frac{1}{G F}\sum_{g=1}^{G}\sum_{f=1}^{F} \lvert \phi_{i,g,c,f} \rvert ,
\end{equation}
\chg{then averaged over windows and normalised across electrodes to give $I_{j,c}$. The resulting cross-fitted attribution magnitude $A_c = \operatorname{median}_j I_{j,c}$ represented the typical importance of electrode $c$, whereas $U_c = \operatorname{median}_j \lvert I_{j,c} - A_c \rvert$ quantified its variation across inner folds. Electrodes were ranked according to}
\begin{equation}
Q_c(\lambda) = A_c - \lambda\, U_c , \qquad \lambda \in \{0, 0.25, 0.5, 1\},
\label{eq:srcfa}
\end{equation}
\chg{where $\lambda = 0$ corresponds to cross-fitted attribution magnitude alone and positive values introduce a penalty for unstable attribution} (the stability-regularised cross-fitted attribution score, SR-CFA). Ties were broken by the lower channel number. The plan did not specify a criterion for choosing $\lambda$, so all four values are reported.

\chg{For each outer fold, each $\lambda$ and each retained electrode count $k \in \{1, \ldots, 16\}$, the classifier was retrained from scratch using only the top-$k$ electrodes and evaluated on the held-out participant.} All models of this electrode-count performance curve used the same fixed, stratified sample of 200,000 training windows of the fold. \chg{This separation between ranking and final evaluation ensured that neither electrode selection nor model training used information from the participant on whom performance was measured.}

\subsection[Hypothesis testing and statistical analysis]{\chg{Hypothesis testing and statistical analysis}}
\label{sec:statistics}

\chg{The primary endpoint was participant-level Macro-F1 across the 16 active gestures;} accuracy and balanced accuracy were also recorded. \chg{For each reduced configuration, performance was compared with the 16-electrode model of the same curve, trained using the same outer-fold data. The paired difference for participant $p$ and subset size $k$ was defined as}
\begin{equation}
\Delta_{p,k} = \mathrm{F1}_{p,k} - \mathrm{F1}_{p,16}.
\end{equation}
\chg{A subset was considered non-inferior (H1) when $L_k$, the lower bound of the 95\% participant-cluster-bootstrap confidence interval for the mean paired difference $\bar\Delta_k$, exceeded the prespecified margin $-\delta$ with $\delta = 0.02$ Macro-F1. The smallest subset satisfying this criterion defined the minimum acceptable electrode count,} $k^* = \min\{k < 16 : L_k > -\delta\}$, recorded as `none' if no $k$ qualifies. As required by the plan, the family of 15 tests was also corrected by the Holm procedure, using one-sided cluster-bootstrap $p$-values for $H_0\colon \bar\Delta_k \le -\delta$ at $\alpha = 0.025$. The margin is an absolute difference in Macro-F1 on a 0--1 scale, fixed before the curve was computed as the largest loss we would accept in exchange for fewer electrodes. Relative to the full-array reference it corresponds to about \MarginPctOfRef\% of performance and about \MarginPctAboveChance\% of the distance above chance (1/16). It is an engineering tolerance and has not been validated against clinical outcomes or user performance. \chg{Because $k^*$ depends on $\delta$, sensitivity analyses were additionally performed at the prespecified margins of 0.01 and 0.03.} Non-inferiority within this tolerance does not imply equivalence.

\chg{To determine whether the attribution ranking conveyed useful information beyond chance (H2), attribution-selected subsets were compared with random subsets of identical size at $k = 4$, 6, 8 and 12. Additional exploratory comparisons included permutation importance, mutual information and geometric farthest-point selection.} Every ranking was computed on the training participants only:
\begin{itemize}
\item SR-CFA ($\lambda = 0.5$) and cross-fitted SHAP ($\lambda = 0$);
\item permutation importance: a model was fitted on 20,000 training windows and each of the 160 features was permuted in 8,000 further windows (3 repeats, Macro-F1 scoring); an electrode's score is the mean over its 10 features;
\item mutual information between each feature and the gesture label on the same 20,000 windows, averaged over the electrode's 10 features;
\item geometric selection: farthest-point sampling on the $2 \times 8$ montage (arc distance around a 24~cm circumference, 2~cm between rings, starting from F1, ties to the lower index), which uses no data;
\item random subsets: 100 independent random subsets per participant at each $k$, and 500 at the prespecified $k^*$ (seed fixed per participant).
\end{itemize}
Evaluating hundreds of random subsets per participant required a lighter comparison model (XGBoost with 50 trees, trained on a stratified sample of 150,000 windows), applied identically to every method; its absolute scores are lower than those of the performance curve and are compared only with each other. The confirmatory H2 family contrasts each attribution ranking with the participant's mean over its random subsets (Holm over 8 tests); comparisons with permutation importance, mutual information and geometric selection are exploratory and uncorrected. We also report the percentile of the SR-CFA subset within each participant's random subsets. Below, `SR-CFA' refers to $\lambda = 0.5$ and `cross-fitted SHAP' to $\lambda = 0$.

\chg{Stability regularisation (H3) was assessed by comparing each $\lambda > 0$ with $\lambda = 0$ on the performance curve at $k = 4$, 6, 8 and 12} (Holm over 12 tests). \chg{Attribution faithfulness (H4) was assessed through deletion with retraining \citep{samek2017evaluating,hooker2019benchmark}: performance was measured after preferentially removing highly ranked (Remove-Top), randomly selected (Remove-Random) or weakly ranked (Remove-Bottom) electrodes.} For $n \in \{2, 4, \ldots, 14\}$ electrodes retained, the comparison model was retrained after removing the $16 - n$ highest- or lowest-ranked electrodes; Remove-Random used the random subsets above and is available for $n \in \{4, 6, 8, 12\}$. A faithful ranking gives Remove-Top $<$ Remove-Random $<$ Remove-Bottom. Paired contrasts were Holm-corrected within H4.

\chg{The participant, rather than the individual signal window, was treated as the statistical unit throughout} ($n = 43$ for GRABMyo, 10 for DB5); no participant was excluded. \chg{Confidence intervals (CI) were obtained as 95\% percentile intervals from 10,000 participant-level cluster-bootstrap resamples} (seed 20260728). Because LOSO training sets overlap, these intervals quantify variation in the observed out-of-fold participant scores and do not fully propagate uncertainty from retraining on a new cohort. \chg{Paired comparisons used the Wilcoxon signed-rank test with the matched-pairs rank-biserial correlation $r_{rb}$ as the effect size, and Holm correction was applied within each confirmatory test family. A non-significant difference was interpreted only as an absence of statistically detectable evidence and not as evidence of equivalence,} because no equivalence test was performed.

\subsection[Exploratory and sensitivity analyses]{\chg{Exploratory and sensitivity analyses}}
\label{sec:exploratory}

\chg{Exploratory analyses examined the stability of electrode rankings across participants, replication of the reduction procedure on NinaPro DB5, cross-session generalisation within GRABMyo, and bipolar re-referencing of the original monopolar recordings. These analyses were interpreted separately from the prespecified primary non-inferiority test.}

\paragraph{Ranking stability} We compared rankings derived from single participants' attributions (all 903 pairs), from the five inner folds and from the 43 outer folds, using Spearman's $\rho$, Kendall's $\tau$ and the Jaccard index of the top-$k$ sets. A participant's attribution profile is the mean of their normalised attribution over the outer folds in which they were held out of an inner model. Outer folds share about 98\% of their training participants, so their agreement is not an independent estimate of between-person stability.

\paragraph{Replication} The LOSO baseline, cross-fitted ranking ($K = 5$) and performance curve were repeated on DB5, with its full array as reference ($k = 16$), using all training windows and the same bootstrap, margins and $k^*$ rule. Holm-corrected results are also reported. No setting was tuned on DB5.

\paragraph{Cross-session} XGBoost was trained on session~1 of $R_{-p}$ and tested on each session of $p$, and trained on sessions 1 and 2 and tested on session~3. These comparisons use the full array and do not test a reduced layout under session transfer. The second condition also has about twice the training data, so it confounds session diversity with data quantity.

\paragraph{Bipolar re-referencing} To use a bipolar pair as the selection unit, we formed $B_j = F_j - F_{j+8}$ ($j = 1, \ldots, 8$) from the monopolar recordings, re-extracted features with the threshold given by the same rule ($4 \times 10^{-3}$~mV; the bipolar signal is about half the monopolar amplitude), and repeated the analysis with the $k = 8$ model as reference. These are derived pairs, not physical bipolar electrodes.

\paragraph{Participant-level aggregation (post hoc)} \chg{A post-hoc extension estimated attribution dispersion directly across participants rather than across pooled inner-fold attribution vectors.} We retained each inner-held-out participant's normalised attribution vector and calculated its median and unscaled median absolute deviation across all outer-training participants. The same score and all four penalties were evaluated at every subset size under the original performance-curve model and training budget, and no penalty was selected from outer-test performance. Exact duplicate electrode sets reused the same fitted evaluation. Comparisons against the original aggregation were paired by participant, and the family of all reduced-size comparisons across penalties received Holm correction. Person vectors share inner models, so their dispersion measures heterogeneity rather than a standard error.

\chg{Departures from the original statistical analysis plan are documented in supplementary methods~S1. These included post-hoc specification of the detailed H2--H4 test families, application of the plan's Holm correction for H1 after the primary analysis, selection of computational training budgets using two folds, and addition of fixed-layout, feature-redundancy and participant-level attribution analyses;} two exploratory analyses listed in the plan were not run. \chg{None altered the prespecified non-inferiority margin, statistical unit, bootstrap procedure or definition of the minimum acceptable electrode count.}
